%%%%%%%%%%%%%%%%%%%%%%%%%%%%%%%%%%%%%%%%%%%%%%%%%%%%%%%%%%%%%%%%%%%%%%
% Lớp 10 | Hệ thức lượng trong tam giác | Giải tam giác và ứng dụng thực tế 
% Dạng 0H4?3-1: Giải tam giác
% Số câu: 21  (sinh tự động bởi tach_cau.py; từ nay thầy sửa trực tiếp file này)
%%%%%%%%%%%%%%%%%%%%%%%%%%%%%%%%%%%%%%%%%%%%%%%%%%%%%%%%%%%%%%%%%%%%%%

%%%%%%%%%%%%%%%% TRẮC NGHIỆM 4 PHƯƠNG ÁN %%%%%%%%%%%%%%%%

% nguồn: MaHoaNH_TracNghiem_Lop10
\begin{ex}%[0H4N3-1]%[ma:b37a9b5b]
	Cho tam giác $ABC$ có diện tích bằng $24$ và chu vi bằng $24$. Bán kính $r$ của đường tròn nội  tiếp tam giác $ABC$ là
	\choice 
		{$r=\dfrac{1}{2}$} 
		{$r=1$} 
		 {$r=3$}
		{\True $r=2$} 
	\loigiai{ 
		Ta có $p=\dfrac{24}{2}=12$.\\ 
	Ta  có $S=p\cdot r\Rightarrow r=\dfrac{S}{p}=2$.}
\end{ex}

% nguồn: MaHoaNH_TracNghiem_Lop10
\begin{ex}%[0H4H3-1]%[ma:8dd06fb9]
	Cho tam giác  $ABC$, biết bán kính đường tròn ngoại tiếp tam giác là $R=15$, diện tích tam giác $ABC$ bằng $45$ và $\widehat{A}=30^\circ$. Tính khoảng cách từ đỉnh $A$ đến cạnh $BC$.
	 \choice {\True $15$}{$20$}{$6$}{$\dfrac{19}{2}$}
	 \loigiai{
		Gọi $d$ là khoảng cách từ đỉnh $A$ đến  cạnh $BC$.\\
		Áp dụng định  lý sinta có $a=2R\sin A=2\cdot 15\cdot\sin30^\circ=15$.\\
		 Mặt khác $S_{ABC}=\dfrac{1}{2}ad$.\\
		Suy ra  $d=\dfrac{2S_{ABC}}{a}=\dfrac{2\cdot45}{15}=6$.\\
		Vậy khoảng cách từ đỉnh $A$ đến cạnh $BC$  là $d=6$.
	 }
\end{ex}

% nguồn: MaHoaNH_TracNghiem_Lop10
\begin{ex}%[0H4Y3-1]%[ma:b0fa8d66]
	Cho đường elip có  phương trình chính tắc $(E)\colon\dfrac{x^2}{4}+\dfrac{y^2}{3}=1$. Một tiêu điểm của elip đó là
	\choice 
		 {$F_2(0; \sqrt{3})$}
		{$F_2(\sqrt{3};  0)$}
		 {\True $F_2(1; 0)$}
		{$F_2(2; 0)$} 
	 \loigiai{
		Ta có $ a^2=4 $, $ b^2=3 $, suy ra $ c=\sqrt{a^2-b^2}=1 $. Do đó tiêu điểm của elip là $ F_{1,2}=\left(\pm 1;0\right)  $.
	} 
\end{ex}

% nguồn: MaHoaNH_TracNghiem_Lop10
\begin{ex}%[0H4Y3-1]%[ma:578c1e32]
	Elip $(E)\colon\dfrac{x^2}{16}+\dfrac{y^2}{9}=1$ có độ  dài trục lớn bằng
	 \choice
		{$9$} 
		{$3$} 
		{\True $8$} 
		 {$4$}
	 \loigiai{
		Với $a^2=16$ ta có $a=4$, do đó độ dài trục lớn  của $(E)$ là $2a=8$.
	} 
\end{ex}

% nguồn: MaHoaNH_TracNghiem_Lop10
\begin{ex}%[0H4B3-1]%[ma:8e2527b0]
	Trong mặt phẳng  tọa độ $Oxy$, cho elip $(E) \colon \dfrac{x^2}{25}+\dfrac{y^2}{9}=1$. Tiêu cự của $(E)$ bằng
	\choice 
		 {$10$}
		{$16$} 
		 {$4$}
		 {$8$}
	\loigiai{ 
		Ta có $a^2=25$, $b^2=9$ suy ra $c=\sqrt{25-9}=4$. Vậy tiêu  cự của elip là $2c=8$.
	 }
\end{ex}

% nguồn: MaHoaNH_TracNghiem_Lop10
\begin{ex}%[0H4Y3-1]%[ma:3b081c6f]
	Trong mặt phẳng $Oxy$,  cho phương trình elip $(E)\colon \dfrac{x^2}{17}+\dfrac{y^2}{9}=1$. Tiêu cự của $(E)$ bằng
	\choice 
		 {\True $4\sqrt{2}$}
		{$\sqrt{26}$} 
		 {$2\sqrt{26}$}
		{$2\sqrt{2}$} 
	\loigiai{ 
		Từ phương trình elip  ta có $a=\sqrt{17}; b=3$.\\
		Vây tiêu cự  là $2c=2\sqrt{a^2-b^2}=4\sqrt{2}$.
	 }
\end{ex}

% nguồn: MaHoaNH_TracNghiem_Lop10
\begin{ex}%[0H4B3-1]%[ma:a65f3da9]
	 Trong mặt phẳng $Oxy$, cho elip $(E)\colon \dfrac{x^2}{25}+\dfrac{y^2}{9}=1$. Khi đó, tiêu cự của elip bằng
	 \choice
		 {10}
		 {6}
		{\True 8} 
		 {4}
	\loigiai{ 
		 Ta có $(E)\colon \dfrac{x^2}{25}+\dfrac{y^2}{9}=1\Rightarrow \heva{&a=5\\&b=3.}$\\
		Tiêu cự của elip  là $2c=\sqrt{a^2-b^2}=2\sqrt{5^2-3^2}=2\cdot4=8$.
	 }
\end{ex}

% nguồn: MaHoaNH_TracNghiem_Lop10
\begin{ex}%[0H4B3-1]%[ma:d37715a6]
	Trong mặt phẳng tọa  độ $Oxy$, cho Elip $(E)\colon\dfrac{x^2}{169}+\dfrac{y^2}{144}=1$. Tiêu cự của $(E)$ là
	\choice 
		{$F_1F_2=26$} 
		{$F_1F_2=24$} 
		{\True  $F_1F_2=10$}
		 {$F_1F_2=5$}
	\loigiai{ 
		Ta có $\heva{&a^2=169\\&b^2=144}\Rightarrow  c^2=a^2-b^2=169-144=25\Leftrightarrow c=5$.\\
		Vậy tiêu cự  của $(E)$ là $F_1F_2=2c=2\cdot 5=10$.
	} 
\end{ex}

% nguồn: MaHoaNH_TracNghiem_Lop10
\begin{ex}%[0H4B3-1]%[ma:b60df049]
	Trong  mặt phẳng tọa độ $Oxy$, cho Elip $(E)\colon\dfrac{x^2}{9}+\dfrac{y^2}{4}=1$. Độ dài trục lớn của $(E)$ là
	\choice 
		{\True  $6$}
		{$3$} 
		{$2$} 
		{$4$} 
	\loigiai{ 
		Ta có  $a^2=9\Rightarrow a=\pm 3$.\\
		Độ dài  trục lớn của $(E)$ là $2a=2\cdot 3=6$.
	} 
\end{ex}

% nguồn: MaHoaNH_TracNghiem_Lop10
\begin{ex}%[0H4B3-1]%[ma:4cac7f29]
	 Trong mặt phẳng tọa độ $Oxy$, cho Elip có phương trình chính tắc $\dfrac{x^2}{34}+\dfrac{y^2}{9}=1$. Điểm nào dưới đây là một tiêu điểm của Elip đã cho?
	 \choice
		{$F_1(-3;0)$} 
		 {\True $F_2(-5;0)$}
		{$F_3(0;-5)$} 
		 {$F_4(\sqrt{34};0)$}
	\loigiai{ 
		 Ta có $a^2=34$, $b^2=9$ suy ra $c^2=a^2-b^2=25$.\\
		Vậy tiêu điểm của Elip đã cho lần  lượt là $(-5;0)$ và $(5;0)$.
	} 
\end{ex}

% nguồn: MaHoaNH_TracNghiem_Lop10
\begin{ex}%[0H4B3-1]%[ma:adebcd58]
	Trong mặt phẳng $Oxy$, cho elip $(E)\colon\dfrac{x^2}{25}+\dfrac{y^2}{9}=1$. Độ  dài tiêu cự của $(E)$ là
	\choice 
		{$4$} 
		{\True $8$} 
		{$16$} 
		{$2$} 
	 \loigiai{
		Elip $(E)\colon\dfrac{x^2}{25}+\dfrac{y^2}{9}=1$ có $\heva{&a^2=25\\&b^2=9}$, do đó $c^2=a^2-b^2=16$.\\ 
		 Vậy tiêu cự của $(E)$ là $2c=2\cdot \sqrt{16}=0$.
	 }
\end{ex}

% nguồn: MaHoaNH_TracNghiem_Lop10
\begin{ex}%[0H4B3-1]%[ma:9f1f05ae]
	Trong mặt phẳng  $Oxy$, cho elip $(E)$ có phương trình chính tắc $\dfrac{x^2}{100}+\dfrac{y^2}{64}=1$. Tiêu cự của elip bằng
	\choice 
		{$16$} 
		 {$20$}
		{\True  $12$}
		{$4\sqrt{41}$} 
	 \loigiai{
		Ta có $\heva{&a^2=100\Rightarrow a=10\\&b^2=64\Rightarrow b=8.}$\\ 
		$\Rightarrow  c=\sqrt{a^2-b^2}=\sqrt{10^2-8^2}=6$.\\
		Tiêu cự là  $2c=12$.
	 }
\end{ex}

% nguồn: MaHoaNH_TracNghiem_Lop10
\begin{ex}%[0H4K3-1]%[ma:17991d2b]
	Cho elip $(E)\colon\dfrac{x^2}{25}+\dfrac{y^2}{9}=1$. Trong các khẳng  định sau, khẳng định nào \textbf{sai}?
	\choice 
		{\True $(E)$  có độ dài trục nhỏ bằng $3$}
		{$(E)$ có các tiêu điểm  $F_1(-4;0)$ và $F_2(4;0)$}
		{$(E)$  có đỉnh $A_1(-5;0)$}
		{$(E)$ có tỉ số $\dfrac{c}{a}=\dfrac{4}{5}$} 
	\loigiai{ 
		Elip $(E)\colon\dfrac{x^2}{25}+\dfrac{y^2}{9}=1$  có $\heva{&a^2=25\\&b^2=9} \Leftrightarrow \heva{&a=5\\&b=3}$, suy ra $c=\sqrt{a^2-b^2}=4$.\\
		Do đó $(E)$ có độ dài  trục nhỏ $2b=6$; $A_1(-5;0)$ là một đỉnh; hai tiêu điểm $F_1(-4;0)$, $F_2(4;0)$ và tỉ số $\dfrac{c}{a}=\dfrac{4}{5}$.\\
		Vậy khẳng định ``$(E)$  có độ dài trục nhỏ bằng $3$'' là sai.
	} 
\end{ex}

% nguồn: MaHoaNH_TracNghiem_Lop10
\begin{ex}%[0H4K3-1]%[ma:cd6b3f10]
	\immini{Khoảng cách từ nhà An ở vị trí $A$ đến nhà Bình ở vị trí $B$ là $500$ m. Từ nhà Bình đi $x$ mét theo phương tạo với $AB$ một góc $60^\circ$ thì sẽ đến nhà bác Mai ở vị trí $M$  và nếu đi thêm $300$ m nữa thì sẽ đến siêu thị ở vị trí $S$ (hình vẽ). Biết rằng quãng đường từ nhà An đến siêu thị bằng $\dfrac{7}{5}$ quãng đường từ nhà An đến nhà bác Mai. Khi đó quãng đường từ nhà Bình đến nhà bác Mai là
	\choice 
		 {\True $625$ m}
		 {$550$ m}
		 {$600$ m}
		{$400$ m}}{\begin{tikzpicture}[scale=0.7]
		 \path
			(0,0)  coordinate (B)
			++(0:6)  coordinate (A)
			(B)++(60:4.5)  coordinate (M)
			(B)++(60:7)  coordinate (S);
		\draw (S)--(B)--(A)--cycle (M)--(A); 
		\foreach \x/\g in  {M/180,B/180,A/0,S/90}
		\fill[black] (\x)  circle (1pt)
			 ($(\g:3mm)+(\x)$) node {$\x$};
		 \path
			(M)--(B)  node[pos=0.5,left]{$x$}
			(M)--(S) node[pos=0.5,left]{$300$} 
			(M)--(A) node[pos=0.5,below]{} 
			 (B)--(A) node[pos=0.5,below]{$500$};
		 \draw
			pic["$60^\circ$",angle radius=15mm]{angle=A--B--S} 
			pic[draw,angle radius=6mm]{angle=A--B--S}; 
	\end{tikzpicture}
		 }
	 \loigiai{
		Vì $x$ là khoảng cách $BM$  nên $x>0$, $BS = x+300$. \\
		Áp dụng định lí  côsin cho tam giác $ABM$:
		$$AM^2 = AB^2 +BM^2 - 2\cdot AB\cdot  BM\cdot \cos 60^\circ = 500^2 +x^2 - 2\cdot 500\cdot x\cdot \dfrac{1}{2}= x^2 -500x +250000.$$
		 Áp dụng định lí cô-sin cho $\triangle ABS$\\
		$AS^2 = AB^2 +BS^2 - 2\cdot AB\cdot BS\cdot \cos 60^\circ = 500^2 +(x+300)^2 -  2\cdot 500\cdot (x+300)\cdot \dfrac{1}{2}\\= x^2 +100x +190000.$\\
		Ta  có:\\
		$AS = \dfrac{7}{5} AM \Leftrightarrow AS^2 = \dfrac{49}{25}AM^2 \Leftrightarrow  25\cdot AS^2 = 49\cdot AM^2$\\
		$\Leftrightarrow 25\cdot  (x^2 +100x +190000) = 49\cdot (x^2 -500x +250000)$\\
		$\Leftrightarrow  24x^2 -27000x + 7500000 =0 \Leftrightarrow \hoac{&x=625 \\ & x=500.}$
	} 
\end{ex}


%%%%%%%%%%%%%%%% ĐÚNG/SAI %%%%%%%%%%%%%%%%

% nguồn: MaHoaNH_DungSai_Lop10
\begin{ex}%[0H4H3-1]%[ma:ceb2239c]
	Cho tam  giác ABC có $BC=a$; $AC=b$; $AB=c$. Biết canh $a=5$ cm; góc $\widehat{B}=83^{\circ}$; $\widehat{C}=57^{\circ}$. Xét tính đúng sai của các khẳng đinh sau
	\choiceTF 
		{\True Số đo góc  $\widehat{A}$ bằng $40^{\circ}$}
		{\True Đô dài $b \approx 7{,}7$ cm (làm tròn môt  chữ số thâp phân)}
		 {Đô dài $c \approx 6{,}3$ cm (tròn môt chữ số thâp phân)}
		{\True Độ dài bán kính đường tròn ngoai tiếp tam giác $R\approx  3{,}9$ cm (làm tròn môt chữ số thập phân)}
	\loigiai{ 
		\begin{itemchoice}
			 \itemch Đúng: $\widehat{A}=180^{\circ}-\widehat{B}-\widehat{C}=180^{\circ}-83^{\circ}-57^{\circ}=40^{\circ}$.
			\itemch Đúng: ta có $\dfrac{a}{\sin A}=\dfrac{b}{\sin B} \Rightarrow \dfrac{5}{\sin 40^{\circ}}=\dfrac{b}{\sin 83^{\circ}} \Rightarrow  b=\dfrac{5\cdot \sin 83^{\circ}}{\sin 40^{\circ}} \approx 7{,}7$.
			\itemch Sai:  $\dfrac{a}{\sin A}=\dfrac{c}{\sin C} \Rightarrow \dfrac{5}{\sin 40^{\circ}}=\dfrac{c}{\sin 57^{\circ}} \Rightarrow c=\dfrac{5\cdot \sin 57^{\circ}}{\sin 40^{\circ}} \approx 6{,}52$.
			\itemch Đúng: $\dfrac{a}{\sin A}=2R\Rightarrow  R=\dfrac{5}{2\sin 40^{\circ}} \approx 3{,}9$.
		\end{itemchoice}
	 }
\end{ex}

% nguồn: MaHoaNH_DungSai_Lop10
\begin{ex}%[0H4V3-1]%[ma:120bc2ad]
	Tam giác $ABC$ có $AB=2$, $AC=1$ và  $\widehat{A}=60^{\circ}$.
	\choiceTF 
		{\True  Độ dài cạnh $BC=\sqrt{3}$}
		{Diện tích tam  giác $ABC$ bằng $\sqrt{3}$}
		 {Độ dài đường cao kẻ từ đỉnh $B$ của tam giác là $h_b=\dfrac{\sqrt{3}}{2}$}
		{\True Bán  kính đường tròn ngoại tiếp tam giác $ABC$ là $R=1$}
	\loigiai{ 
		\begin{itemchoice}
			\itemch Ta có $BC^2=AB^2+AC^2-2\cdot AB\cdot AC\cos\widehat{A}=2^2+1^2-2\cdot 2\cdot 1\cos 60^{\circ}=3$ nên  $BC=\sqrt{3}$.
			\itemch  Ta có $S_{\triangle ABC}=\dfrac{1}{2}\cdot AB\cdot AC\cdot\sin A=\dfrac{\sqrt{3}}{2}$.
			\itemch Do  $h_b=\dfrac{2S_{\triangle ABC}}{b}=\dfrac{2\cdot\dfrac{\sqrt{3}}{2}}{1}=\sqrt{3}$.
			\itemch Ta có  $\dfrac{BC}{\sin A}=2R\Leftrightarrow\dfrac{\sqrt{3}}{\sin 60^{\circ}}=2R\Leftrightarrow R=1$.
		\end{itemchoice}
	} 
\end{ex}

% nguồn: MaHoaNH_DungSai_Lop10
\begin{ex}%[0H4V3-1]%[ma:21a0c25e]
	Cho tam  giác $ABC$ có $BC=12\sqrt{3}$; $A=120^{\circ}$; $B=45^{\circ}$.
	\choiceTF 
		{\True Độ dài  cạnh $AC=12\sqrt{2}$ và $AB=6(\sqrt{6}-\sqrt{2})$}
		{Diện tích tam giác $ABC$ bằng  $54-18\sqrt{3}$}
		{Độ dài đường trung tuyến kẻ từ đỉnh  $A$ của tam giác là $m_a=36(5-2\sqrt{3})$}
		 {\True Bán kính đường tròn nội tiếp tam giác $ABC$ là $r=9\sqrt{2}+3\sqrt{6}-18$}
	 \loigiai{
		\begin{itemchoice}
			\itemch  Xét tam giác $ABC$ có
			$\widehat{A}+\widehat{B}+\widehat{C}=180^{\circ}$ (định lí  tổng $3$ góc trong tam giác).\\
			Suy ra $\widehat{C}=180^{\circ}-\widehat{A}-\widehat{B}=180^{\circ}-120^{\circ}-45^{\circ}=15^{\circ}  $.\\
			Áp  dụng định lí sintrong tam giác $ABC$, ta có\\
			$\dfrac{BC}{\sin A}=\dfrac{AC}{\sin  B}\Rightarrow AC=\dfrac{BC\cdot\sin B}{\sin A} =\dfrac{12\sqrt{3}\cdot\sin 45^{\circ}}{\sin 120^{\circ}}=12\sqrt{2}$.\\
			 Vậy $AC=12\sqrt{2}$.\\
			Mà $\dfrac{BC}{\sin A}=\dfrac{AB}{\sin C}\Rightarrow AB=\dfrac{BC\cdot\sin  C}{\sin A} =\dfrac{12\sqrt{3}\cdot\sin 15^{\circ}}{\sin 120^{\circ}}=6\left(\sqrt{6}-\sqrt{2}\right)$.\\
			Vậy  $AB=6\left(\sqrt{6}-\sqrt{2}\right)$.
			\itemch Diện tích tam  giác $ABC$ là\\
			$S_{ABC}=\dfrac{1}{2}AC\cdot BC\cdot\sin C =\dfrac{1}{2}\cdot  12\sqrt{2}\cdot 12\sqrt{3}\cdot\sin 15^{\circ}=108-36\sqrt{3}$.\\
			 Vậy $S_{ABC}=108-36\sqrt{3}$.
			\itemch Do  $m_a^2=\dfrac{b^2+c^2}{2}-\dfrac{a^2}{4}=\dfrac{(12\sqrt{2})^2+36\left(\sqrt{6}-\sqrt{2}\right)^2}{2}-\dfrac{(12\sqrt{3})^2}{4}=36\left(5-2\sqrt{3}\right)$.\\ Suy ra $m_a=6\sqrt{5-2\sqrt{3}}$.
			\itemch Ta có $S=p\cdot  r\Rightarrow r=\dfrac{S}{p}=2\cdot\dfrac{108-36\sqrt{3}}{12\sqrt{3}+12\sqrt{2}+6\sqrt{6}-6\sqrt{2}}=9\sqrt{2}+3\sqrt{6}-18$.
		\end{itemchoice}
	 }
\end{ex}

% nguồn: MaHoaNH_DungSai_Lop10
\begin{ex}%[0H4V3-1]%[ma:25ea2bf7]
	Cho tam giác $ABC$ có $a=13, b=14$, $c=15$. Khi  đó:
	 \choiceTF
		{\True Tam giác $ABC$ có  $\cos A=0{,}6$}
		{Chu vi tam giác  $ABC$ là $21$}
		{\True Diện tích của tam giác $ABC$  bằng $84$}
		{Đường tròn ngoại tiếp tam giác  $ABC$ có chu vi bằng $32{,}5\pi$}
	 \loigiai{
		\begin{itemchoice}
			 \itemch Đúng.\\
			Áp  dụng định lý côsinta có $\cos A=\dfrac{b^2+c^2-a^2}{2bc}=\dfrac{14^2+15^2-13^2}{2\cdot 14\cdot 15}=0{,}6$.
			\itemch  Sai.\\
			Chu vi tam giác  $ABC$ là $CV=a+b+c=13+14+15=42$.
			 \itemch Đúng.\\
			 Nửa chu vi tam giác $ABC$ là $p=\dfrac{a+b+c}{2}=\dfrac{13+14+15}{2}=21$.
			 Diện tích của tam giác $ABC$ là
			\[S=\sqrt{p\cdot  (p-a)\cdot (p-b)\cdot (p-c)}=\sqrt{21\cdot (21-13)\cdot (21-14)\cdot (21-15)}=84.\]
			 \itemch Sai.\\
			Gọi $R$ là bán kính đường tròn ngoại tiếp tam giác  $ABC$.\\
			Ta có $S=\dfrac{abc}{4R} \Rightarrow R=\dfrac{abc}{4S}=\dfrac{13\cdot 14\cdot 15}{4\cdot 84}=\dfrac{65}{8}$.\\ 
			Chu  vi đường tròn ngoại tiếp tam giác $ABC$ là
			\[CV=2\pi  \dfrac{65}{8}=\dfrac{65\pi}{4}=16{,}25\pi.
			 \]
		\end{itemchoice}
	 }
\end{ex}


%%%%%%%%%%%%%%%% TỰ LUẬN %%%%%%%%%%%%%%%%

% nguồn: MaHoaNH_TuLuan_Lop10
\begin{bt}%[0H4H3-1]%[ma:5359b064]
	Cho tam giác $ABC$,  biết $AB=3$; $AC=4$; $\widehat{A}=150^{\circ}$. Tìm diện tích của tam giác $ABC$
	 \loigiai{
		Với $AB=c$;  $AC=b$; $BC=a$.\\
		Ta có $S_{ABC}=\dfrac{1}{2} bc\sin A=\dfrac{1}{2} 3\cdot 4\cdot  \sin 150^{\circ}=3$.
	} 
\end{bt}

% nguồn: MaHoaNH_TuLuan_Lop10
\begin{bt}%[0H4V3-1]%[ma:1bcc0509]
	Cho tam giác $ABC$ có $\widehat{B}=60^{\circ}$, $BC=8$, $AB+AC=12$. Tính độ dài  các cạnh $AB$, $AC$.
	\loigiai{ 
		Đặt $AB=x$ ($x>0$). Khi  đó, $AC=12-x$ ($x<12$).\\
		Áp dụng định lí côsin cho cạnh  $AC$:\\
		$AC^{2}=AB^{2}+BC^{2}-2AB\cdot BC\cdot\cos B$\\ 
		$(12-x)^{2}=x^{2}+8^{2}-2\cdot x\cdot8\cdot\cos  60^{\circ}$\\
		 $144-24x+x^{2}=x^{2}+64-16x\cdot\dfrac{1}{2}$\\
		$144-24x+x^{2}=x^{2}+64-8x$\\ 
		$144-64=24x-8x \Leftrightarrow 80=16x  \Rightarrow x=5$.\\
		Vậy độ dài cạnh $AB=5$, độ  dài cạnh $AC=12-5=7$.
	} 
\end{bt}

% nguồn: MaHoaNH_TuLuan_Lop10
\begin{bt}%[0H4K3-1]%[ma:44682627]
	Để cắt một biển quảng cáo hình Elip có trục lớn là $80$ cm và trục nhỏ là $40$ cm, từ một tấm Inox hình chữ nhật có kích thước $80\times 40$, người ta ghim hai chiếc đinh tại hai tiêu điểm của Elip và vẽ hình  Elip đó lên tấm Inox (như hình vẽ minh họa dưới đây). Hỏi phải ghim hai chiếc đinh cách các mép tấm Inox bao nhiêu và lấy vòng dây không đàn hồi có độ dài bao nhiêu?
	\begin{center}
		\begin{tikzpicture}[line cap=butt,line join=miter,>=stealth]
				 \tikzset{declare function={xmin=-3; xmax=3; ymin=-2; ymax=2;},
					smooth,samples=450 
				 }
				 \path (0,0) coordinate (O);
				\draw[->] (xmin,0)--(xmax,0)  node[shift={(-100:7pt)},font=\normalsize]{$ x $};
				 \draw[->] (0,ymin)--(0,ymax) node[shift={(190:7pt)},font=\normalsize]{$ y $};
				\fill (0,0) node[shift={(225:7pt)},font=\normalsize]{$  O $};
				\clip (xmin,ymin)  rectangle (xmax,ymax);
				\fill (-2,0) circle (1pt); \fill (2,0) circle (1pt);  \fill (0,-1) circle (1pt); \fill (0,1) circle (1pt); \fill (-1.2,0) circle (1pt); \fill (1.2,0) circle (1pt);
				 \draw (O) ellipse (2 and 1)
					(-2,1) node[left]{$P$}--(2,1) node[right]{$Q$}--(2,-1) node[right]{$R$}--(-2,-1)  node[left]{$S$}--(-2,1)
					(0,1) node[above left]{$B_2$} 
					(0,1)  node[above right]{$b$}
					 (0,-1) node[below left]{$B_1$}
					(0,-1) node[below  right]{$-b$}
					(-2,0) node[above  left]{$-a$}
					 (-2,0) node[below left]{$A_1$}
					 (2,0) node[above right]{$a$}
					 (2,0) node[below right]{$A_2$}
					(-1.2,0) node[above]{$F_1$} 
					 (1.2,0) node[above]{$F_2$};
			\end{tikzpicture}
	\end{center}
	\loigiai{ 
		Ghim đinh cách mép  tấm Inox: $A_1F_1=A_2F_2=a-c=40-20\sqrt{3}$.\\
		 Độ dài vòng dây không đàn hồi: $2a+2c=80+40\sqrt{2}$.
	 }
\end{bt}
